\section{案例研究与实验设计}

\subsection{机器人示例：Pick-and-place pipeline}

\begin{figure}[H]
\centering
\includegraphics[width=0.8\textwidth]{slide-02-03.jpg}
\caption{Slide 02-03 展示了一条 pick-and-place pipeline，从演示数据、生成模型到 rollout monitoring。}
\end{figure}

该 pipeline 内部将 expert trajectory 聚合、训练 diffusion policy，并在 scoring server 中实时计算 KL drift。部署前 post-mortem 包含: 1) control length 2) intervention log 3) visual replay diff。

\begin{table}[H]
\centering
\begin{tabular}{p{0.25\textwidth} p{0.7\textwidth}}
\toprule
阶段 & 核心测量 \\
\midrule
数据 & Coverage of pick/place, human override rate, sensor noise \\
训练 & Loss plateau, latent interpolation, sampling temperature drift \\
部署 & KL drift, rejection rate, human override latency \\
\bottomrule
\end{tabular}
\caption{Pick-and-place 模仿学习 pipeline 的监控点}
\end{table}

\begin{knowledgebox}{录像回放的作用}
把 policy rollouts 录成 short clips，与 expert clip 做 diff，有助于 PM 快速识别 drift 来源，从而决定是否 rollback。
\end{knowledgebox}

\subsection{Evaluation Matrix 与对齐实验}

\begin{figure}[H]
\centering
\includegraphics[width=0.8\textwidth]{slide-02-07.jpg}
\caption{Slide 02-07 把 benchmark、safety 与 governance 放在一个分层矩阵中，便于 cross-team 对齐。}
\end{figure}

Evaluation matrix 由三层组成：benchmark scores、safety metrics、ops guardrails。每个版本都需通过如下 tests：
\begin{itemize}
    \item Benchmark：MMLU-like tasks、LongBench multi-hop。
    \item Safety：forbidden content rejection、hallucination detectors。
    \item Ops：latency、human override、drift ticket count。
\end{itemize}

\begin{warningbox}{过拟合 benchmark 的风险}
盲目追 benchmark 会让 policy 忽略 rare prompt；一定要把 safety 和 ops guardrail 作为 gating 条件，防止“指标驱动”的 drift。
\end{warningbox}

\subsection{本章小结}

案例研究把演示数据、生成策略、evaluation matrix 和治理转成具体的 artifact（dashboard、table、clips），为后续扩展提供可复制的实验框架。

